\section{Introduction}
\label{sec:intro}

% Android malware detection takes an Android Package Kit (APK) as input and outputs a probability of the APK being malicious.
% Machine learning techniques are widely adopted to assist the analysis of suspicious APKs and summarization of malicious patterns, improving the detection scalability and accuracy~\cite{qiu2020survey,liu2022deep,wu2019malscan,wu2021homdroid,mclaughlin2017deep,he2022msdroid,xu2018deeprefiner,aafer2013droidapiminer,arp2014drebin}.
% Traditional methods involve manual analysis of suspicious APKs and summarization of malicious patterns as hand-crafted rules~\cite{enck2009lightweight,felt2011android}, but it is not scalable or effective.
% % It falls short in handling the surfeit of new Android software produced annually and proves ineffective against previously unknown malware.
% Instead, recent advancements leverage machine learning (ML) to improve the scalability and accuracy in Android malware detection~\cite{qiu2020survey,liu2022deep,wu2019malscan,wu2021homdroid,mclaughlin2017deep,he2022msdroid,xu2018deeprefiner,aafer2013droidapiminer,arp2014drebin}.
% Specifically, the common paradigm is extracting features from APKs (\textit{e.g.,} permissions and intents), encoding them into numeric vectors, and applying ML models (\textit{e.g.,}neural networks) to automatically distinguish malware from benign apps.

Over the past decade, ML-based Android malware detection has received increased attention from various research communities, such as software engineering, security, and machine learning~\cite{qiu2020survey,liu2022deep,wu2019malscan,wu2021homdroid,mclaughlin2017deep,he2022msdroid,xu2018deeprefiner,aafer2013droidapiminer,arp2014drebin}.
Much attention has been given to exploring various combinations of APK features and ML models.
The trend is primarily led by advancements in ML models (\textit{e.g.}, graph neural network~\cite{he2022msdroid}) and analogies drawn from other well-studied fields (\textit{e.g.}, social network~\cite{wu2019malscan}).
Most existing approaches report high F1 scores (up to 0.98)~\cite{wu2019malscan,kim2018multimodal}.
Such promising results motivate us to ask a number of important research questions.
For example,
how does the existing work represent and incorporate diverse APK features into ML models for Android malware detection?
How do different detection approaches compare when evaluated on the same datasets, metrics, and toolchains?
Are current ML-based approaches sufficient to meet the requirements in real-world scenarios?
Does a more powerful ML model necessarily yield better results?
Does incorporating more features to describe app behaviors always come with better outcomes?
Is there a positive correlation between the detection effectiveness and efficiency?

In this paper, we investigate the current state and offer an in-depth understanding of ML-based Android malware detection with empirical and quantitative analysis.
Given the significant growth of app availability --- Google Play's app count has reached 3.95 millions, up 6.5\% from 2023~\cite{appnumber} --- this study focuses on assessing methods that are scalable to large datasets and practical in real-world usage. 
We aim to experimentally compare the results from a number of representative approaches to understand the advancements and challenges encountered in the field.
While instructive, we identify several challenges that hinder the comparison.
\begin{itemize}[leftmargin=10pt, itemsep=2pt]
% \begin{description}%[leftmargin=0pt]
  %\item \textit{Existing experimental results are not directly comparable.}
  \item \textit{Unfair Comparisons.}
  Previous approaches are usually evaluated on datasets of different sizes, goodware-to-malware ratios, and training-to-testing ratios~\cite{arp2014drebin,hou2017hindroid,li2021robust}.
  Moreover, they report outcomes using diverse metrics (\textit{e.g.,} F1-score, Accuracy, and False Positive Rate), making it unclear which method performs better under specific settings.
  In addition, they often rely on different toolchains to develop detectors, which introduce ambiguity --- whether the promising results stem from the novelty of the approach or not~\cite{vasilache2018tensor}.

  %\item[Real-world scenarios, \textit{e.g.,} malware evolution, obfuscation, and adversarial attacks are not often fully considered in many existing studies.]
  \item \textit{Unrealistic Evaluations.} 
  Android malware detectors face a rapid threat landscape~\cite{barbero2022transcending}, with malware continually evolving to evade detection.
  The growing usage of obfuscation also makes malware harder to detect as its malicious intentions are hidden~\cite{aonzo2020obfuscapk}.
  Additionally, ML models can be tricked by intentionally perturbed inputs~\cite{pierazzi2020problemspace,carlini2017towards,suya2020hybrid,grosse2017adversarial}.
  Although the impacts of some of these scenarios have been studied~\cite{jordaney2017transcend, barbero2022transcending,pendlebury2019tesseract,ceschin2020machine}, the lack of a comprehensive assessment creates a gap in our understanding of the main challenges that hinder the adoption of Android malware detection in real-world scenarios.

  %\item[Most methods overlook reporting the end-to-end efficiency from feature engineering to ML modeling.]
  \item \textit{Unclear Computational Costs.}
  With the exponential growth of apps in sizes and complexities, feature extraction gradually becomes notably time-consuming. 
  Furthermore, as ML models evolve in complexity, they demand greater computational power to achieve state-of-the-art results.
  Unfortunately, it remains uncertain what prices to pay to gain the desired results.
\end{itemize}

%{\noindent \bf Our Approach.}To tackle these issues, we design a general-purpose framework, \codename, to evaluate ML-based Android malware detection solutions. 

To tackle these issues, we design a general-purpose framework, \codename, to evaluate ML-based Android malware detection solutions. 
\codename is modular and configurable, facilitating the development of new Android malware detectors and the design of various evaluation scenarios.
Specifically, we dissect the detection pipeline into three phases: APK characterization, feature representation, and ML modeling.
Guided by this taxonomy, we first empirically analyze the related work to explore how current approaches detect Android malware.
% We begin with dissecting the detection pipeline into three phases: APK characterization, feature representation, and ML modeling.
% Guided by this taxonomy, we empirically analyze the related work to explore how current approaches detect Android malware.
%
Next, we apply our framework to 12 representative approaches from three research communities: software engineering~\cite{wu2019malscan,wu2021homdroid,wu2021android}, security~\cite{arp2014drebin,mariconti2016mamadroid,mclaughlin2017deep,xu2018deeprefiner,kim2018multimodal,xu2020sdac,he2022msdroid}, and machine learning~\cite{hou2017hindroid,li2021robust}.
In particular, we ensure that the same task (\textit{e.g,} feature extraction) across different methods is performed using the same toolchain, eliminating potential discrepancies in the results caused by different toolchains.
To further facilitate the development of new detectors, we implement \codename to be modular and configurable, where components (\textit{e.g.,} neural networks) can be freely swapped.
For evaluation, we randomly sample 221,310 apps spanning ten years from AndroZoo~\cite{allix2016androzoo}, a continuously updated public repository of Android apps.
The malware ratio is set as 10\% to mimic a realistic setting~\cite{pendlebury2019tesseract}.
To make the experiments mirror real-world conditions, we evaluate the selected approaches using common used metrics (\textit{e.g.,} F1-score and Accuracy), focusing on the following aspects: the effectiveness and efficiency in malware detection, the robustness against app evolution and obfuscation, and the resilience against adversarial attacks.

% More importantly, we observe a {\em key aspect} of effectiveness of these solutions: {\em domain semantics}, which is whether the effort is "related" to the problem itself.
% For example,  we discover that powerful ML models are not a silver bullet for better detection outcomes, as ML models may not have the right input to describe the problem. 
% As another example, while APK features are important to depict apps, simply adding more features counterintuitively can lead to negative impacts, as non-related features dilute the relevant domain semantics.
\noindent \textbf{Findings and Recommendations.} Through our empirical and quantitative analysis, we draw a holistic picture of the state of the art in ML-based Android malware detection.
Specifically, APK characterization and ML modeling remain central themes.
Many recent detectors achieve comparable results under the same experimental settings.
However, their effectiveness still falls short in challenging scenarios, such as those with limited data volumes or under adversarial attacks.
% Moreover, we observe a key aspect behind effective solutions --- \textit{eigen semantics}, the unique characteristic meaning of app behaviors and features.
We attribute that the ability of extracting semantic information that describes malicious behaviors from APK features, which we term as \textit{malware semantics}, is the key to the effectiveness of ML-based detectors.
For example, from APK feature aspect, we find that simply adding more features to describe apps can lead to negative impacts, as non-related features dilute the relevant semantic information.
Additionally, from model aspect, we notice that employing powerful ML models does not necessarily yield better detection outcomes, as the models may not have the right input to describe app behaviors.
Another interesting observation is that increased computational overhead is not a reliable indicator of enhanced detection capabilities.
Future work should pay more attention to designing robust and practical detectors for real-world usage, with balancing effectiveness and efficiency in mind.
More discussions are in Sec.~\ref{sec:findings}.

We emphasize that the design of effective ML-based Android malware detection solutions should be guided by the extraction and integration of malware semantics describing app behaviors from APK features.
Defining and quantifying these semantics remains an open research challenge but is critical for advancing ML-based detection methods.
We hope our findings help shape future research directions in the field, including, but not limited to, feature selection, model design, and the efficient allocation of computing resources.

In summary, we make the following contributions:
\begin{itemize}[leftmargin=10pt, topsep=0pt, itemsep=2pt]
  \item
  We conduct a thorough systematic investigation of ML-based Android malware detection using empirical and qualitative methods, drawing a holistic picture of the field.
  \item
  We design a general-purpose framework, \codename, to facilitate the implementation and evaluation of various Android malware detectors.
  For comparison in realistic settings, we collect and open-source the largest dataset to date, both in size and time span.
  The framework, along with the dataset, is available at \url{https://anonymous.4open.science/r/RkRyb2lk}.
  \item 
  We offer a comprehensive comparative analysis of 12 representative approaches using \codename, focusing on assessing their effectiveness, robustness, and efficiency. We point out that the key for enhancing ML-based Android malware detection is incorporating the malware semantics derived from APK features.
\end{itemize}

